% !TEX program = pdflatex
% ------------------------------------------------------------
% MATH 521 Homework Template (adapted from your MATH 421 HW10)
%
% HOW TO USE (quick):
%   1) Edit the Metadata block (HW number / due date especially).
%   2) Paste each problem statement under \problem[]{<n>}.
%   3) Write in the Scratch/Discussion box if you want.
%   4) Write your proof under "Proof." and end with \qed (or \qedhere).
%
% Notes:
% - Keeps theorem/definition boxes (tcolorbox).
% - Keeps ToC + PDF bookmarks for each problem.
% - Adds a Scratch/Discussion box per problem (optional).
% - Header bar has course + meeting info (left) and HW number (right).
% ------------------------------------------------------------

% TROUBLESHOOTING NOTE:
% If you ever see errors like `\tcb@insert@after@part` or other tcolorbox internals being
% an "Undefined control sequence", that usually means LaTeX is accidentally loading a
% DIFFERENT (often older) `tcolorbox.sty` from your project folder or local texmf tree.
% Fix: search your repo for `tcolorbox.sty` and delete/rename any local copy so TeX uses
% the TeX Live version. Also try "Clean" in LaTeX Workshop to remove stale .aux/.toc/.out.

\documentclass[12pt]{article}

% =========================
% 0) Metadata (EDIT ME)
% =========================
\newcommand{\CourseCode}{MATH 521}
\newcommand{\CourseTitle}{Analysis I}
\newcommand{\CourseSection}{LEC 002}
\newcommand{\Semester}{Spring 2026}
\newcommand{\Instructor}{Thierry Laurens}
% \newcommand{\MeetingInfo}{MWF 1:20--2:10 \;|\; VV B119} % optional
\newcommand{\HomeworkNumber}{6} % <-- change each week
\newcommand{\YourName}{Alejandro De La Torre}
\newcommand{\DueDate}{\today} % <-- or set e.g. January 31, 2026

% =========================
% 1) Core math + proof tools
% =========================
\usepackage{amsmath, amssymb, amsthm}

% =========================
% 2) Lists and spacing
% =========================
\usepackage{enumitem}
\setlist[enumerate,1]{label=\textbf{(\alph*)}, leftmargin=2em, itemsep=0.25em}
\setlist[itemize]{leftmargin=1.5em, itemsep=0.25em}
\setlength{\parskip}{0.35em}
\setlength{\parindent}{0pt}

% =========================
% 3) Page geometry + header/footer
% =========================
\usepackage{geometry}
\geometry{margin=1in}

\usepackage{fancyhdr}
\pagestyle{fancy}
\fancyhf{}
\lhead{\CourseCode\ --- \CourseSection, \Semester}
\rhead{Homework \HomeworkNumber}
\cfoot{\thepage}
\renewcommand{\headrulewidth}{0.4pt}
% Fixes common fancyhdr warning about header height:
\setlength{\headheight}{15pt}
% Optional: reclaim a bit of vertical space after increasing headheight
\addtolength{\topmargin}{-2pt}

% =========================
% 4) Links (subtle)
% =========================
\usepackage[hidelinks]{hyperref}
\setcounter{tocdepth}{1}

% =========================
% 5) Quotes / figures
% =========================
\usepackage{csquotes}
\usepackage{graphicx}
\graphicspath{{images/}} % put figures in ./images/

% =========================
% 6) Simple scratch box (no tcolorbox)
% =========================
% A lightweight environment for brief planning / scratch work.
% This avoids any tcolorbox dependencies.
\newenvironment{scratch}{%
  \par\smallskip
  \noindent\textbf{Discussion \& Scratch Work}\begin{quote}\small
}{%
  \end{quote}\smallskip
}

% =========================
% 7) Lightweight theorem-like environments (optional)
% =========================
\newtheorem*{claim}{Claim}
\newtheorem*{lemma}{Lemma}
\newtheorem*{remark}{Remark}
\newtheorem{theorem}{Theorem}
\newtheorem{proposition}{Proposition}
\newtheorem{corollary}{Corollary}
\newtheorem{definition}{Definition}
\newtheorem{example}{Example}
\newtheorem{exercise}{Exercise}

% =========================
% 8) Problem header command (with TOC + PDF bookmarks)
% Usage: \problem[Optional short title]{<number>}
% =========================
\makeatletter
\newcommand{\problem}[2][]{%
  \def\prob@title{Problem #2\if\relax\detokenize{#1}\relax\else\ (\textit{#1})\fi}%
  \phantomsection%
  \pdfbookmark[1]{\prob@title}{prob#2}%
  \addcontentsline{toc}{section}{\prob@title}%
  \section*{\prob@title}%
  \vspace{-0.25em}\hrule\vspace{0.8em}%
}
\makeatother

% =========================
% 9) Title block
% =========================
\title{Homework \HomeworkNumber}
\author{\YourName \\
\CourseCode\ --- \CourseTitle \\
\CourseSection \\
Instructor: \Instructor}
\date{\DueDate}

\begin{document}
\maketitle

% ------------------ collaboration + sources ------------------
% \section*{Collaboration Statement}
% Write “I worked alone” or list collaborators / external resources.
% "I worked on this homework independently. No outside collaborators or resources were used."

\tableofcontents
\bigskip
\hrule
\bigskip

\newpage

% ============================================================
% Problems 
% ============================================================

\problem[]{1}
Prove that the series
\[
\sum_{n=1}^{\infty} \frac{1}{(3n-2)(3n+1)}
\]
converges, and find its value. (Hint: You computed the partial sums on Homework 1.)

\begin{scratch}
We want to show the series converges and compute its sum. The hint suggests
computing the partial sums and looking for a telescoping form. A standard move
is to use partial fractions:
\[
\frac{1}{(3n-2)(3n+1)} = \frac{1}{3}\left(\frac{1}{3n-2}-\frac{1}{3n+1}\right),
\]
which should make the partial sums collapse.
\end{scratch}

\textbf{Proof.}\\
Let
\[
a_n := \frac{1}{(3n-2)(3n+1)} \qquad (n\in\mathbb{N}).
\]
For $N\in\mathbb{N}$ define the $N$th partial sum
\[
S_N := \sum_{n=1}^{N} a_n.
\]
Using partial fractions,
\[
\frac{1}{(3n-2)(3n+1)}
= \frac{1}{3}\left(\frac{1}{3n-2}-\frac{1}{3n+1}\right).
\]
Hence
\[
S_N
= \frac{1}{3}\sum_{n=1}^{N}\left(\frac{1}{3n-2}-\frac{1}{3n+1}\right)
= \frac{1}{3}\Bigg[\left(1-\frac{1}{4}\right)+\left(\frac{1}{4}-\frac{1}{7}\right)+\cdots+\left(\frac{1}{3N-2}-\frac{1}{3N+1}\right)\Bigg].
\]
All intermediate terms cancel, so
\[
S_N = \frac{1}{3}\left(1-\frac{1}{3N+1}\right)=\frac{N}{3N+1}.
\]
Therefore
\[
\lim_{N\to\infty} S_N = \lim_{N\to\infty}\frac{N}{3N+1}=\frac{1}{3}.
\]
Since the sequence of partial sums $(S_N)$ converges, the series
$\sum_{n=1}^{\infty} \frac{1}{(3n-2)(3n+1)}$ converges, and its value is $\frac{1}{3}$.
\qed

\newpage

\problem[]{2}
Let $\sum_{n=1}^{\infty} a_n$ and $\sum_{n=1}^{\infty} b_n$ be convergent series.
\begin{enumerate}[label=\textbf{(\alph*)}]
\item Prove that for any $k\in\mathbb{R}$, the series $\sum_{n=1}^{\infty} k a_n$
converges and
\[
\sum_{n=1}^{\infty} k a_n = k \sum_{n=1}^{\infty} a_n.
\]
\item Prove that the series $\sum_{n=1}^{\infty} (a_n+b_n)$ converges and
\[
\sum_{n=1}^{\infty} (a_n+b_n) = \sum_{n=1}^{\infty} a_n + \sum_{n=1}^{\infty} b_n.
\]
\end{enumerate}

\begin{scratch}
\end{scratch}

\textbf{Proof.}\\
\qed

\newpage

\problem[]{3}
\begin{enumerate}[label=\textbf{(\alph*)}]
\item Prove that
\[
\sum_{n=1}^{\infty} \frac{n-1}{2^{n+1}}=\frac{1}{2}.
\]
(Hint: $\frac{k-1}{2^{k+1}}=\frac{k}{2^{k}}-\frac{k+1}{2^{k+1}}$.)
\item Use part (a) to compute
\[
\sum_{n=1}^{\infty} \frac{n}{2^{n}}.
\]
\end{enumerate}

\begin{scratch}
\end{scratch}

\textbf{Proof.}\\
\qed

\newpage

\problem[]{4}
Prove that if $\sum_{n=1}^{\infty} a_n$ converges absolutely and
$\{b_n\}_{n\ge 1}$ is a bounded sequence, then the series
$\sum_{n=1}^{\infty} a_n b_n$ also converges. (Hint: Use the Cauchy criterion.)

\begin{scratch}
\end{scratch}

\textbf{Proof.}\\
\qed

\newpage

\problem[]{5}
\begin{enumerate}[label=\textbf{(\alph*)}]
\item Prove that if $\sum_{n=1}^{\infty} a_n$ converges absolutely, then the
series $\sum_{n=1}^{\infty} a_n^2$ also converges.
\item Provide an example where $\sum_{n=1}^{\infty} a_n$ diverges and
$\sum_{n=1}^{\infty} a_n^2$ converges.
\end{enumerate}

\begin{scratch}
\end{scratch}

\textbf{Proof.}\\
\qed

\newpage

\problem[]{6}
For each of the following series, determine if it converges and prove your
answer.
\begin{enumerate}[label=\textbf{(\alph*)}]
\item \[
\sum_{n=2}^{\infty} \frac{1}{[n+(-1)^n]^2}
\]
\item \[
\sum_{n=1}^{\infty} \left(\sqrt{n+1}-\sqrt{n}\right)
\]
\item \[
\sum_{n=1}^{\infty} (-1)^n
\]
\end{enumerate}

\begin{scratch}
\end{scratch}

\textbf{Proof.}\\
\qed

% ============================================================
% Appendix
% ============================================================
\newpage
\appendix
\section*{Appendix: Toolbox}
\addcontentsline{toc}{section}{Appendix: Toolbox}
\subsection*{Ces\`aro--Stolz}
\textbf{Theorem (Ces\`aro--Stolz).} Let $\{a_n\}_{n\ge 1}$ be a sequence of
nonzero real numbers. Then
\[
\liminf_{n\to\infty}\left|\frac{a_{n+1}}{a_n}\right|
\le \liminf_{n\to\infty} |a_n|^{1/n}
\le \limsup_{n\to\infty} |a_n|^{1/n}
\le \limsup_{n\to\infty}\left|\frac{a_{n+1}}{a_n}\right|.
\]

\textbf{Lemma.} If $r>0$, then $\lim_{n\to\infty} r^{1/n}=1$.

\subsection*{Series: Definitions and Basic Facts}
\textbf{Definition (Series).} Let $\{a_n\}_{n\ge 1}$ be a sequence of real
numbers. The series $\sum_{n=1}^{\infty} a_n$ converges/diverges if the sequence
of partial sums $s_n=a_1+\cdots+a_n$ converges/diverges.

\textbf{Definition (Absolute Convergence).} The series $\sum_{n=1}^{\infty} a_n$
converges absolutely if the series $\sum_{n=1}^{\infty} |a_n|$ converges.

\textbf{Remark.} As $|a_n|\ge 0$, the sequence
$s_n=|a_1|+\cdots+|a_n|$ is increasing, so either
$\sum_{n=1}^{\infty} |a_n|$ converges or diverges to $+\infty$.

\textbf{Corollary.} If $\sum_{n=1}^{\infty} a_n$ converges, then
$\lim_{n\to\infty} a_n=0$.

\textbf{Remark (Geometric Series).} The series $\sum_{n=0}^{\infty} r^n$
converges if $|r|<1$ and
\[
\sum_{n=0}^{\infty} r^n=\frac{1}{1-r}.
\]
If $|r|\ge 1$, then $\sum_{n=0}^{\infty} r^n$ diverges.

\subsection*{Cauchy Criterion for Series}
\textbf{Theorem (Cauchy criterion).} The series $\sum_{n=1}^{\infty} a_n$
converges if and only if for all $\varepsilon>0$ there exists $N\in\mathbb{N}$
such that
\[
\left|\sum_{k=m}^{n} a_k\right|<\varepsilon
\quad\text{for all } n\ge m\ge N.
\]

\subsection*{Absolute Convergence}
\textbf{Corollary.} If $\sum_{n=1}^{\infty} a_n$ converges absolutely, then it
converges.

\subsection*{Comparison Test}
\textbf{Theorem (Comparison test).} Let $\sum_{n=1}^{\infty} a_n$ be a series
with $a_n\ge 0$ for all $n\ge 1$.
\begin{enumerate}[label=\textbf{(\arabic*)}]
\item If $\sum_{n=1}^{\infty} a_n$ converges and $|b_n|\le a_n$ for all $n\ge 1$,
then $\sum_{n=1}^{\infty} b_n$ converges.
\item If $\sum_{n=1}^{\infty} a_n$ diverges and $b_n\ge a_n$ for all $n\ge 1$,
then $\sum_{n=1}^{\infty} b_n$ diverges.
\end{enumerate}

\subsection*{Dyadic Criterion}
\textbf{Theorem (Dyadic criterion).} Let $\{a_n\}_{n\ge 1}$ be a decreasing
sequence with $a_n\ge 0$ for all $n\ge 1$. The series
$\sum_{n=1}^{\infty} a_n$ converges if and only if the series
$\sum_{n=0}^{\infty} 2^n a_{2^n}$ converges.

\textbf{Corollary.} The series $\sum_{n=1}^{\infty} \frac{1}{n^{\alpha}}$
converges if and only if $\alpha>1$.
\end{document}
